% ---------------------------------------------------------------------
% Figura: la pipeline adottata in questo lavoro.
%
% Richiede nel preambolo di tesi.tex:
%   \usepackage{tikz}
%   \usetikzlibrary{arrows.meta,positioning}
%
% Si inserisce con % ---------------------------------------------------------------------
% Figura: la pipeline adottata in questo lavoro.
%
% Richiede nel preambolo di tesi.tex:
%   \usepackage{tikz}
%   \usetikzlibrary{arrows.meta,positioning}
%
% Si inserisce con % ---------------------------------------------------------------------
% Figura: la pipeline adottata in questo lavoro.
%
% Richiede nel preambolo di tesi.tex:
%   \usepackage{tikz}
%   \usetikzlibrary{arrows.meta,positioning}
%
% Si inserisce con % ---------------------------------------------------------------------
% Figura: la pipeline adottata in questo lavoro.
%
% Richiede nel preambolo di tesi.tex:
%   \usepackage{tikz}
%   \usetikzlibrary{arrows.meta,positioning}
%
% Si inserisce con \input{fig_pipeline} dove serve.
% Struttura ispirata alla Figura 1 di Hoang et al. (IEEE Access, 2023),
% ma con un solo ramo percorso: il ramo appreso e' tratteggiato perche'
% e' stato confrontato nel Capitolo 4 e non schierato sul dispositivo.
% ---------------------------------------------------------------------

\begin{figure}[htbp]
\centering
\begin{tikzpicture}[
  font=\footnotesize,
  blocco/.style   = {draw, rounded corners=2pt, align=center,
                     text width=20mm, minimum height=10mm, inner sep=2pt},
  dato/.style     = {blocco, fill=black!4},
  scelto/.style   = {blocco, fill=black!10, very thick},
  alt/.style      = {blocco, dashed, draw=black!45, text=black!55},
  fr/.style       = {-{Stealth[length=1.8mm]}, thick},
  fralt/.style    = {-{Stealth[length=1.8mm]}, dashed, draw=black!45}
]

% --- riga 1: dal video alla finestra temporale ---
\node[dato] (video) at (0,0)    {Video RGB};
\node[dato] (posa)  at (2.7,0)  {MediaPipe Pose\\33 landmark};
\node[dato] (feat)  at (5.4,0)  {Feature\\geometriche\\e integrali};
\node[dato] (dtw)   at (8.1,0)  {Selezione DTW\\12 su 33};
\node[dato] (fin)   at (10.8,0) {Finestre\\$L=22$, passo 2};

\draw[fr] (video) -- (posa);
\draw[fr] (posa)  -- (feat);
\draw[fr] (feat)  -- (dtw);
\draw[fr] (dtw)   -- (fin);

% --- riga 2: dalla finestra all'allarme (verso sinistra) ---
\node[scelto] (rf)      at (10.8,-2.3) {Random Forest};
\node[alt]    (reti)    at (10.8,-4.1) {GRU, LSTM,\\CNN 1D,\\Transformer,\\grafo};
\node[dato]   (soglia)  at (8.1,-2.3)  {Soglia $0{,}50$};
\node[dato]   (regola)  at (5.4,-2.3)  {Cinque finestre\\consecutive};
\node[dato]   (allarme) at (2.7,-2.3)  {Allarme};

\draw[fr] (fin)    -- (rf);
\draw[fr] (rf)     -- (soglia);
\draw[fr] (soglia) -- (regola);
\draw[fr] (regola) -- (allarme);

% ramo confrontato ma non schierato
\draw[fralt] (fin.east)  -- ++(0.7,0) |- (reti.east);
\draw[fralt] (reti.west) -| (soglia.south);

% --- protocollo di valutazione ---
\node[align=center, text width=34mm, draw, dotted, rounded corners=2pt,
      inner sep=3pt] (val) at (2.0,-4.1)
      {Valutazione\\LOSO su 10 soggetti\\$F_2$};
\draw[fralt, draw=black!60] (allarme.south) -- (val.north);

\end{tikzpicture}

\caption{La catena adottata in questo lavoro, dal fotogramma all'allarme. Il
ramo continuo e' la configurazione installata sul dispositivo; il ramo
tratteggiato raccoglie le architetture confrontate nel
Capitolo~\ref{cap:protocollo} a parita' di rappresentazione, finestre e
protocollo, ma non schierate. La struttura segue quella generale degli
approcci skeleton-based descritta in~\cite{HO}, percorsa pero' in un solo
ramo.}
\label{fig:pipeline}
\end{figure}
 dove serve.
% Struttura ispirata alla Figura 1 di Hoang et al. (IEEE Access, 2023),
% ma con un solo ramo percorso: il ramo appreso e' tratteggiato perche'
% e' stato confrontato nel Capitolo 4 e non schierato sul dispositivo.
% ---------------------------------------------------------------------

\begin{figure}[htbp]
\centering
\begin{tikzpicture}[
  font=\footnotesize,
  blocco/.style   = {draw, rounded corners=2pt, align=center,
                     text width=20mm, minimum height=10mm, inner sep=2pt},
  dato/.style     = {blocco, fill=black!4},
  scelto/.style   = {blocco, fill=black!10, very thick},
  alt/.style      = {blocco, dashed, draw=black!45, text=black!55},
  fr/.style       = {-{Stealth[length=1.8mm]}, thick},
  fralt/.style    = {-{Stealth[length=1.8mm]}, dashed, draw=black!45}
]

% --- riga 1: dal video alla finestra temporale ---
\node[dato] (video) at (0,0)    {Video RGB};
\node[dato] (posa)  at (2.7,0)  {MediaPipe Pose\\33 landmark};
\node[dato] (feat)  at (5.4,0)  {Feature\\geometriche\\e integrali};
\node[dato] (dtw)   at (8.1,0)  {Selezione DTW\\12 su 33};
\node[dato] (fin)   at (10.8,0) {Finestre\\$L=22$, passo 2};

\draw[fr] (video) -- (posa);
\draw[fr] (posa)  -- (feat);
\draw[fr] (feat)  -- (dtw);
\draw[fr] (dtw)   -- (fin);

% --- riga 2: dalla finestra all'allarme (verso sinistra) ---
\node[scelto] (rf)      at (10.8,-2.3) {Random Forest};
\node[alt]    (reti)    at (10.8,-4.1) {GRU, LSTM,\\CNN 1D,\\Transformer,\\grafo};
\node[dato]   (soglia)  at (8.1,-2.3)  {Soglia $0{,}50$};
\node[dato]   (regola)  at (5.4,-2.3)  {Cinque finestre\\consecutive};
\node[dato]   (allarme) at (2.7,-2.3)  {Allarme};

\draw[fr] (fin)    -- (rf);
\draw[fr] (rf)     -- (soglia);
\draw[fr] (soglia) -- (regola);
\draw[fr] (regola) -- (allarme);

% ramo confrontato ma non schierato
\draw[fralt] (fin.east)  -- ++(0.7,0) |- (reti.east);
\draw[fralt] (reti.west) -| (soglia.south);

% --- protocollo di valutazione ---
\node[align=center, text width=34mm, draw, dotted, rounded corners=2pt,
      inner sep=3pt] (val) at (2.0,-4.1)
      {Valutazione\\LOSO su 10 soggetti\\$F_2$};
\draw[fralt, draw=black!60] (allarme.south) -- (val.north);

\end{tikzpicture}

\caption{La catena adottata in questo lavoro, dal fotogramma all'allarme. Il
ramo continuo e' la configurazione installata sul dispositivo; il ramo
tratteggiato raccoglie le architetture confrontate nel
Capitolo~\ref{cap:protocollo} a parita' di rappresentazione, finestre e
protocollo, ma non schierate. La struttura segue quella generale degli
approcci skeleton-based descritta in~\cite{HO}, percorsa pero' in un solo
ramo.}
\label{fig:pipeline}
\end{figure}
 dove serve.
% Struttura ispirata alla Figura 1 di Hoang et al. (IEEE Access, 2023),
% ma con un solo ramo percorso: il ramo appreso e' tratteggiato perche'
% e' stato confrontato nel Capitolo 4 e non schierato sul dispositivo.
% ---------------------------------------------------------------------

\begin{figure}[htbp]
\centering
\begin{tikzpicture}[
  font=\footnotesize,
  blocco/.style   = {draw, rounded corners=2pt, align=center,
                     text width=20mm, minimum height=10mm, inner sep=2pt},
  dato/.style     = {blocco, fill=black!4},
  scelto/.style   = {blocco, fill=black!10, very thick},
  alt/.style      = {blocco, dashed, draw=black!45, text=black!55},
  fr/.style       = {-{Stealth[length=1.8mm]}, thick},
  fralt/.style    = {-{Stealth[length=1.8mm]}, dashed, draw=black!45}
]

% --- riga 1: dal video alla finestra temporale ---
\node[dato] (video) at (0,0)    {Video RGB};
\node[dato] (posa)  at (2.7,0)  {MediaPipe Pose\\33 landmark};
\node[dato] (feat)  at (5.4,0)  {Feature\\geometriche\\e integrali};
\node[dato] (dtw)   at (8.1,0)  {Selezione DTW\\12 su 33};
\node[dato] (fin)   at (10.8,0) {Finestre\\$L=22$, passo 2};

\draw[fr] (video) -- (posa);
\draw[fr] (posa)  -- (feat);
\draw[fr] (feat)  -- (dtw);
\draw[fr] (dtw)   -- (fin);

% --- riga 2: dalla finestra all'allarme (verso sinistra) ---
\node[scelto] (rf)      at (10.8,-2.3) {Random Forest};
\node[alt]    (reti)    at (10.8,-4.1) {GRU, LSTM,\\CNN 1D,\\Transformer,\\grafo};
\node[dato]   (soglia)  at (8.1,-2.3)  {Soglia $0{,}50$};
\node[dato]   (regola)  at (5.4,-2.3)  {Cinque finestre\\consecutive};
\node[dato]   (allarme) at (2.7,-2.3)  {Allarme};

\draw[fr] (fin)    -- (rf);
\draw[fr] (rf)     -- (soglia);
\draw[fr] (soglia) -- (regola);
\draw[fr] (regola) -- (allarme);

% ramo confrontato ma non schierato
\draw[fralt] (fin.east)  -- ++(0.7,0) |- (reti.east);
\draw[fralt] (reti.west) -| (soglia.south);

% --- protocollo di valutazione ---
\node[align=center, text width=34mm, draw, dotted, rounded corners=2pt,
      inner sep=3pt] (val) at (2.0,-4.1)
      {Valutazione\\LOSO su 10 soggetti\\$F_2$};
\draw[fralt, draw=black!60] (allarme.south) -- (val.north);

\end{tikzpicture}

\caption{La catena adottata in questo lavoro, dal fotogramma all'allarme. Il
ramo continuo e' la configurazione installata sul dispositivo; il ramo
tratteggiato raccoglie le architetture confrontate nel
Capitolo~\ref{cap:protocollo} a parita' di rappresentazione, finestre e
protocollo, ma non schierate. La struttura segue quella generale degli
approcci skeleton-based descritta in~\cite{HO}, percorsa pero' in un solo
ramo.}
\label{fig:pipeline}
\end{figure}
 dove serve.
% Struttura ispirata alla Figura 1 di Hoang et al. (IEEE Access, 2023),
% ma con un solo ramo percorso: il ramo appreso e' tratteggiato perche'
% e' stato confrontato nel Capitolo 4 e non schierato sul dispositivo.
% ---------------------------------------------------------------------

\begin{figure}[htbp]
\centering
\begin{tikzpicture}[
  font=\footnotesize,
  blocco/.style   = {draw, rounded corners=2pt, align=center,
                     text width=20mm, minimum height=10mm, inner sep=2pt},
  dato/.style     = {blocco, fill=black!4},
  scelto/.style   = {blocco, fill=black!10, very thick},
  alt/.style      = {blocco, dashed, draw=black!45, text=black!55},
  fr/.style       = {-{Stealth[length=1.8mm]}, thick},
  fralt/.style    = {-{Stealth[length=1.8mm]}, dashed, draw=black!45}
]

% --- riga 1: dal video alla finestra temporale ---
\node[dato] (video) at (0,0)    {Video RGB};
\node[dato] (posa)  at (2.7,0)  {MediaPipe Pose\\33 landmark};
\node[dato] (feat)  at (5.4,0)  {Feature\\geometriche\\e integrali};
\node[dato] (dtw)   at (8.1,0)  {Selezione DTW\\12 su 33};
\node[dato] (fin)   at (10.8,0) {Finestre\\$L=22$, passo 2};

\draw[fr] (video) -- (posa);
\draw[fr] (posa)  -- (feat);
\draw[fr] (feat)  -- (dtw);
\draw[fr] (dtw)   -- (fin);

% --- riga 2: dalla finestra all'allarme (verso sinistra) ---
\node[scelto] (rf)      at (10.8,-2.3) {Random Forest};
\node[alt]    (reti)    at (10.8,-4.1) {GRU, LSTM,\\CNN 1D,\\Transformer,\\grafo};
\node[dato]   (soglia)  at (8.1,-2.3)  {Soglia $0{,}50$};
\node[dato]   (regola)  at (5.4,-2.3)  {Cinque finestre\\consecutive};
\node[dato]   (allarme) at (2.7,-2.3)  {Allarme};

\draw[fr] (fin)    -- (rf);
\draw[fr] (rf)     -- (soglia);
\draw[fr] (soglia) -- (regola);
\draw[fr] (regola) -- (allarme);

% ramo confrontato ma non schierato
\draw[fralt] (fin.east)  -- ++(0.7,0) |- (reti.east);
\draw[fralt] (reti.west) -| (soglia.south);

% --- protocollo di valutazione ---
\node[align=center, text width=34mm, draw, dotted, rounded corners=2pt,
      inner sep=3pt] (val) at (2.0,-4.1)
      {Valutazione\\LOSO su 10 soggetti\\$F_2$};
\draw[fralt, draw=black!60] (allarme.south) -- (val.north);

\end{tikzpicture}

\caption{La catena adottata in questo lavoro, dal fotogramma all'allarme. Il
ramo continuo e' la configurazione installata sul dispositivo; il ramo
tratteggiato raccoglie le architetture confrontate nel
Capitolo~\ref{cap:protocollo} a parita' di rappresentazione, finestre e
protocollo, ma non schierate. La struttura segue quella generale degli
approcci skeleton-based descritta in~\cite{HO}, percorsa pero' in un solo
ramo.}
\label{fig:pipeline}
\end{figure}
